\section{Exact examples and proof diagnostics}
\label{sec:examples}

Three small examples separate the notions used above before we report the
bounded falsification controls.

\begin{example}[Prime base]
\label{ex:prime-three}
  Let $p=3$.  Then
  \[
    r_1=1,\qquad r_2=4,\qquad r_3=13,
  \]
  so the first hole classes are $1+3\Z$, $4+9\Z$, and $13+27\Z$.
  They are exact skeleton holes for every frozen directive.  The initial
  block of length $3^N$ has essential period $3^{N+1}$; in particular, the
  three coordinates of $B_1$ share period $9$ and no smaller positive
  common period.
\end{example}

\begin{example}[Composite base]
\label{ex:composite-four}
  Let $p=4$.  The exact skeleton periods are still $4^N$, with hole
  residues $r_1=1$, $r_2=5$, and $r_3=21$.  However, choosing the prime
  divisor $\ell=2$ in \cref{thm:constructive-split} gives the common period
  \[
    2\cdot4^N<4^{N+1}
  \]
  for every initial block $B_N$.  At $N=1$, the strict counterperiod is
  $8<16$.  This example prevents one from conflating an essential skeleton
  period with the common period of a finite collection of positions.
\end{example}

\begin{example}[Least-period reduction after quotienting]
\label{ex:period-reduction}
  The directive $(0,1,0,2)^{\Z}$ is frozen.  Merging the nonadjacent letters
  $1$ and $2$ is admissible and produces $(0,a,0,a)^{\Z}$, whose least
  period is two rather than four.  This is why quotient directives are
  normalized after applying the block map.  By contrast, merging $0$ with
  $1$ identifies an adjacency edge and leaves the frozen target family.
\end{example}

\subsection{Finite falsification controls}
\label{subsec:diagnostics}

The proof above is deductive and independent of computation.  We
nevertheless used deterministic finite searches to catch indexing errors,
improper composite-base inversions, wrong refinement directions, and
unlicensed scope extensions.  \Cref{tab:diagnostics} reports selected
counts from three frozen checkers.  The Stage-2 implementation compares the
direct affine formula with a nested-hole evaluator and includes a bounded
local-rule census.  A reciprocal checker independently implements its
arithmetic, directive, partition, and chromatic routines.  A third checker
uses larger skeleton, constructiveness, and partition ranges.  None imports
the theorem proof as code.

\begin{table}[t]
\centering
\caption{Deterministic finite falsification controls reproduced from
the frozen candidate and audit receipts.  The implementations use distinct evaluators,
and all typed mutations are rejected.  These bounded counts are not used
to prove any theorem.}
\label{tab:diagnostics}
\small
\begin{tabular}{@{}lrrr@{}}
\toprule
Control & Stage-2 & Reciprocal & Root \\
\midrule
Formula / nested-hole comparisons & 68,288 & 24,024 & -- \\
Skeleton checks & 28,764 & 1,519 & 984,576 \\
Prime smaller periods rejected & 3,918 & 19,766 & 200,988 \\
Composite progression / shift checks & 998,025 & 80,563 & 426,088 \\
High-center identities & 1,920 & 7,200 & 3,744 \\
Partition tests & 477 & 1,632 & 9,874 \\
Bounded local-rule cases & 972 & -- & -- \\
Typed mutation controls & 4 & 4 & -- \\
\midrule
Local-rule false positives / negatives & $0/0$ & -- & -- \\
\bottomrule
\end{tabular}
\par\vspace{2pt}
{\footnotesize Finite falsification controls only.  Lock prefixes:
candidate \texttt{c070bd76d8a2}, reciprocal \texttt{8b4d54a8cc7c}, root \texttt{e1dca456aecf}.}
\end{table}


The local-rule census tested $972$ bounded source/target/radius/base cases;
its $132$ context-consistent cases were exactly its $132$ letter-quotient
cases, with no false positive or false negative.  This equality motivated a
search for a uniform proof but is not used in
\cref{thm:factor-collapse}.  Typed mutations include a wrong-base identity
proposal, a nonpointed shift, an adjacent merge, a nonsurjective target
map, and the false all-base constructive assertion.  Rejection of a
wrong-base or nonpointed mutation means only that it lies outside the frozen
contract, not that no such map exists in a larger category.

The full deterministic protocol, input hashes, and rerun command are given
in \cref{app:reproducibility}.  All integer counts in
\cref{tab:diagnostics} are exact for the declared finite ranges, and every
infinite statement is closed separately in \cref{sec:skeletons,sec:factors,sec:poset}.
